\documentclass[11pt]{article}
\usepackage{amsmath,amssymb,amsthm}
\usepackage[margin=1in]{geometry}

\newtheorem{theorem}{Theorem}
\newtheorem{corollary}{Corollary}

\title{L1 --- Equivalence Completion (Paper 11):\\
Hypercharge Normalization and Electric Charge Quantization\\
under Admissible Redescription}
\author{Amos Jay Maley}
\date{}

\begin{document}
\maketitle

\section*{0. Scope and Non-Negotiables}
This paper is eliminative. It introduces no new primitives, no new ontology,
no new anchors, and no new tensor load.
All results are invariant under admissible redescription.

We work with fixed chiral fermion content,
fixed gauge representations,
and a single electroweak Higgs doublet.
No additional U(1) factors, no exotic charges,
and no compensating sectors are permitted.

\section*{1. Apparent Freedom of Hypercharge Normalization}
In gauge-phase bookkeeping,
hypercharge assignments appear freely scalable
by an overall normalization factor.
Electric charge quantization is then often treated
as an external phenomenological input.

In broken-phase bookkeeping,
electric charges appear sharply quantized
after electroweak symmetry breaking,
suggesting a possible asymmetry between the two skins.

\section*{2. Admissibility Envelope}
Allowed redescriptions include:
\begin{itemize}
\item overall rescalings preserving relative charge assignments,
\item admissible transport between gauge-phase and broken-phase skins,
\item reparameterizations preserving operator existence and gauge invariance.
\end{itemize}

Any normalization choice that fails transport
without importing new structure is inadmissible.

\section*{3. Transport of Charge Data}
Electric charge is defined as the unbroken generator
$Q = T_3 + Y$,
where $T_3$ is fixed by the non-Abelian representation content.

Transport from the gauge-phase skin to the broken-phase skin
requires that hypercharge assignments permit
a single unbroken $U(1)_{\mathrm{em}}$
acting consistently on all fermions and on the Higgs field.

This requirement constrains hypercharge normalization
beyond a mere conventional scaling.

\section*{4. Collapse of Normalization Freedom}
Assume arbitrary hypercharge normalization
independent of the declared fermion content.
Transport to the broken-phase skin.

Generic normalizations fail to yield
integer electric charges
consistent across all fermion multiplets
without introducing:
(i) additional gauge structure,
(ii) exotic fractional representations, or
(iii) compensating fields.
Each option violates the declared load or admissibility constraints.

Therefore, admissible transport collapses
the apparent normalization freedom.

\section*{5. Core Result}
\begin{theorem}[Hypercharge Normalization and Charge Quantization Equivalence]
Fix the chiral fermion content,
gauge representations,
and a single electroweak Higgs doublet.
Under admissible redescription and global constraints
(no new load, no patches),
hypercharge normalization is fixed up to an overall sign,
and electric charge is necessarily quantized.

Charge quantization is not an additional assumption,
but a transport necessity.
\end{theorem}

\begin{corollary}[Uniqueness of Electric Charge Assignment]
All admissible electric charge assignments
are unique up to overall sign and normalization
fixed by the declared fermion content.
No additional continuous freedom exists.
\end{corollary}

\begin{corollary}[Skin Illusion of Hypercharge Freedom]
The apparent freedom to rescale hypercharge arbitrarily
is a bookkeeping artifact of gauge-phase representation.
Cross-skin admissibility collapses this illusion.
\end{corollary}

\section*{6. Closure Statement}
Hypercharge normalization and electric charge quantization
are not independent inputs.
They are enforced jointly by admissible transport
under fixed chiral load and Higgs structure.
No external assumptions are required.

\end{document}

